Most structural systems are designed to remain fixed throughout their service lives. This fixed-form design philosophy has produced reliable structures with well-defined load paths, efficient use of material, and predictable load response. However, emerging applications such as post-disaster replacement infrastructure and compactly stowable systems for space exploration require structures that can change shape and withstand loads during storage, deployment, and operation. Existing deployable and reconfigurable systems offer several mechanisms for shape transformation. However, scaling these mechanisms to large load-bearing structures is challenging, as the complex joints required for shape transformations are prone to localized stress concentrations and buckling under eccentric loading. Furthermore, metamaterials that use elastic deformations to achieve functional multistability rely on specialized materials studied in small-scale architectures, which are often expensive and difficult to integrate into larger structures. The central challenge is to introduce reconfigurability into structures without sacrificing load capacity, stability, or application-specific mechanical response.

To that end, this dissertation introduces a framework to transform static structural topologies into functional, reconfigurable systems. Trusses are chosen to develop this framework due to their simple geometry, well-defined load paths, and high stiffness-to-weight ratio. The framework converts triangular units in static trusses into flat-foldable quadrilateral linkages by adding joints in accordance with Grashof's flat-foldability criterion and member-force information from a prescribed design load. This transformation makes trusses reconfigurable while preserving the load capacity of the original topology. The dissertation also develops a sequential kinematic analysis tool to evaluate deployment trajectories, actuation sequences, and packing efficiency. Fabrication-aware design strategies account for member thicknesses, joint clearances, and physical interference to prevent overlap and interlocking during deployment. Load tests on proof-of-concept prototypes demonstrate that transformed trusses can achieve large geometric reconfiguration while retaining the stiffness and load capacity of static counterparts.

Next, a Three-node Torsional Spring (3NTS) is developed to introduce rotational stiffness at the joints of reconfigurable trusses. The 3NTS provides a buffer against load-reversal-induced buckling and enables functional applications such as controlled deployment, energy storage, and multistability. The 3NTS is validated against analytical solutions to capture elastic buckling instabilities and large deformations. We build on the 3NTS to model nonlinear bi- and multi-stable rotational hinges, including a hinge design that uses magnetic interactions to create multiple stable configurations. We show that reconfigurable trusses augmented with these hinges exhibit sequential collapse and multistability under loading. By tuning the strength of the magnetic interactions, the hinge response can be programmed to support applications such as reusable impact absorption, controlled standoff in military vehicles, and path planning in robotics.

Finally, we extend this dissertation's philosophy beyond civil structures by shifting the focus towards curved-origami-inspired shape-morphing hulls. This work demonstrates how curved-origami principles help create rapidly deployable hulls from flat sheets. These curved-origami hulls can be designed to match conventional hull forms and morph shape to tune hydrodynamic performance on demand.

Together, the work presented in this dissertation provides a framework for integrating isolated reconfigurable mechanisms into existing static structures to obtain structural systems that combine controlled motion with programmable mechanical response.